% Generated draft: requires mathematical and source review.
\documentclass[11pt]{article}
\usepackage[margin=25mm]{geometry}
\usepackage{fontspec}
\setmainfont{Latin Modern Roman}
\usepackage{amsmath,amssymb,mathtools,unicode-math}
\setmathfont{Latin Modern Math}
\usepackage{enumitem,needspace,xurl,hyperref}
\setlength{\parindent}{0pt}
\setlength{\parskip}{5pt}
\setlength{\emergencystretch}{3em}
\newcommand{\N}{\mathbb{N}}
\newcommand{\Z}{\mathbb{Z}}
\newcommand{\Q}{\mathbb{Q}}
\newcommand{\R}{\mathbb{R}}
\newcommand{\C}{\mathbb{C}}
\newcommand{\F}{\mathbb{F}}
\newcommand{\A}{\mathbb{A}}
\newcommand{\PP}{\mathbb{P}}
\newcommand{\E}{\mathbb{E}}
\newcommand{\eps}{\varepsilon}
\newcommand{\dd}{\,\mathrm{d}}
\begin{document}
\section*{UnsolvedMath ID: 30002904}
\textbf{Source problem code:} catalogue entry
\subsection{Definitions and mathematical statement}
Let \(n\) denote the integer matrix size in the source, and put
\[
G=\mathrm{SL}_n(\Z)=\{g\in M_n(\Z):\det g=1\}.
\]
The source does not specify the intended range of \(n\) or its quantification. Let \(\lVert\cdot\rVert\) denote the norm appearing in the source, and, for real \(X>0\), define the finite set
\[
B_X=\{g\in G:\lVert g\rVert\le X\}.
\]
The background describes an ordinary Euclidean-ball model, but the supplied statement does not define the norm precisely.

For \(g_1,g_2\in G\), let \(H=\langle g_1,g_2\rangle\) be the subgroup they generate, and let \([G:H]\) be the number of left cosets of \(H\) in \(G\). Interpreting the sampling as two independent uniform draws from \(B_X\), with replacement, the question is whether
\[
\lim_{X\to\infty}
\frac{\#\{(g_1,g_2)\in B_X\times B_X:
[G:\langle g_1,g_2\rangle]=\infty\}}
{(\#B_X)^2}
=1.
\]
The quotient is considered for sufficiently large \(X\), when \(B_X\ne\varnothing\). Independence is the interpretation used to write this formula; it is not explicitly stated in the supplied record. The intended dimension range and precise Euclidean norm remain unresolved.

The supplied curation labels the problem partially solved and reports generic thin free groups in a symmetrized Euclidean-ball model, while describing this ordinary Euclidean-ball question as not fully settled. This is a supplied literature qualification, not an independently verified result or a replacement of the question above.
\subsection{Short English statement}
Choose two integral determinant-one matrices uniformly from a growing norm ball. Does the probability that their generated subgroup has infinite index tend to one? The displayed formulation interprets the choices as independent. The source leaves the dimension range and precise norm unspecified; its supplied partial-resolution assessment concerns a symmetrized model.
\subsection{Sources}
\begin{itemize}
\item Local UnsolvedMath registry: \texttt{data/unsolvedmath/problems.json}, with \texttt{catalog/problems.json} as the supplementary registry; permanent ID \(30002904\), exact source code \texttt{catalogue entry}. The supplied record belongs to the Oberwolfach Reports Open Problems collection. No individual proposer is identified; the supplied proposed year is 2015.
\item Supplied source citation: \emph{New Perspectives on the Interplay between Discrete Groups in Low-Dimensional Topology and Arithmetic Lattices} (2016), Oberwolfach Reports, pp.\ 1701--1746, DOI 10.4171/owr/2015/30. Original-report link: \url{https://doi.org/10.4171/owr/2015/30}.
\item Supplied literature source, labeled ``Fuchs and Rivin primary paper'': \url{https://academic.oup.com/imrn/article/2017/17/5385/3056825}. No paper title or fuller author information is supplied.
\item Literature assessment checked at \texttt{2026-08-21}: the supplied curation labels the problem partially solved and reports generic thin free groups in a symmetrized Euclidean-ball model, while stating that the ordinary Euclidean-ball question is not fully settled. This assessment is supplied dated curation and has not been independently verified here. The background explicitly requests independent verification of time-sensitive status and cited claims.
\end{itemize}
\end{document}
